\chapter{Literature Survey}
This chapter surveys existing research, industry reports, and commercial tools relevant to non-human identity governance, AI-driven security automation, and policy-as-code enforcement. The review identifies the theoretical foundations underpinning the proposed NHI Governance Agent and highlights persistent gaps in current approaches that motivate this work.

\section[Introduction to the Survey]{\textbf{Introduction to the Survey}}
The literature review spans three domains: academic research on identity governance, anomaly detection, and zero-trust architectures; industry reports and surveys quantifying the scale and severity of the NHI threat landscape; and analysis of existing commercial tools that address parts of the NHI governance problem. Rather than treating each source in isolation, this survey synthesises findings thematically to build a coherent picture of the current state of the art and the opportunities for innovation.

\section[NHI Threat Landscape]{\textbf{NHI Threat Landscape and Scale of the Problem}}
The proliferation of cloud-native architectures, microservice deployments, and DevOps automation pipelines has produced a massive expansion in the number of non-human identities operating within enterprise environments. Enterprise-wide analyses have revealed that \glspl{nhi} now outnumber human identities at ratios ranging from 10:1 to as high as 144:1, with year-over-year growth rates approaching 44\% \cite{Entro2025}. This exponential growth is driven by the increasing reliance on service principals, \gls{api} keys, deploy tokens, managed identities, and \gls{cicd} pipeline secrets that enable machine-to-machine communication at cloud scale.

The security implications of this growth are severe. Credential-based attacks consistently rank among the most prevalent breach vectors globally, with credential abuse identified as the primary attack path in nearly half of all recorded data breaches \cite{VerizonDBIR2023}. The \gls{owasp} Non-Human Identities Top~10 report identifies the most critical risks facing organisations, including improper offboarding of machine credentials, secret leakage in source code repositories, and overprivileged service accounts \cite{OWASP2025}. The same report found that 70\% of organisations have plaintext secrets stored in code repositories and 47\% maintain unused \gls{iam} roles that remain active and exploitable.

Surveys of security practitioners further underscore the severity of the challenge. A survey of over 800 IT and security professionals found that only 15\% of organisations report high confidence in their ability to prevent NHI-related attacks, while 69\% express significant concern about the state of their machine identity security posture \cite{CSA2024}. The primary pain points identified include inadequate auditing capabilities, insufficient privilege management controls, and the absence of automated policy enforcement for non-human credentials. Analyst reports on the NHI management market confirm that the machine-to-human identity ratio exceeds 17:1 even in moderately sized organisations and project significant market growth as enterprises recognise NHI governance as a board-level security priority \cite{KuppingerCole2025, ResearchMarkets2024}. Credential access tactics documented in established threat intelligence frameworks demonstrate that adversaries actively target long-lived, unrotated machine credentials as a reliable means of establishing persistent access to compromised environments \cite{MITREAttack}.

\section[Zero Trust and IAM Frameworks]{\textbf{Zero Trust and Identity Access Management Frameworks}}
The zero trust security model has emerged as the dominant paradigm for securing modern cloud environments, replacing the traditional perimeter-based approach with the principle that no entity---human or machine---should be implicitly trusted. The foundational architecture defines three core tenets: verify explicitly, enforce least privilege, and assume breach \cite{NIST800207}. While originally formulated for human users, these principles apply with equal urgency to non-human identities that operate autonomously and often hold elevated privileges.

Research on \gls{iam} in cloud environments has explored challenges posed by dynamic resource provisioning and ephemeral workloads. Studies examining IAM with a focus on \gls{zta} identify the tension between automated credential provisioning and the requirement for continuous verification \cite{Mungoli2023}. Cross-cloud federated IAM systems using trust score computation and behavioural modelling have been proposed to enforce least-privilege access across cloud boundaries \cite{Potluri2024}. Surveys of IAM in multi-cloud environments document the fragmented state of identity governance practices \cite{CloudIAMSurvey2023}.

The emergence of autonomous AI agents introduces additional governance challenges. Research on zero-trust identity frameworks for agentic AI demonstrates the inadequacy of conventional authentication protocols for autonomous agents \cite{Narajala2025}. Complementary work on decentralised identity management integrates self-sovereign identity principles and zero-knowledge proofs into multi-domain orchestration frameworks \cite{Bernabe2025}. Digital identity guidelines and international standards for information security management provide the compliance frameworks against which NHI governance solutions must be evaluated \cite{NIST80063, ISO27001, SOC2AICPA}.

\section[AI-Driven Anomaly Detection]{\textbf{AI-Driven Trust Assessment and Anomaly Detection}}
The application of machine learning to security operations has gained significant traction as organisations seek to move beyond static, rule-based monitoring. Comprehensive surveys of anomaly detection methodologies categorise approaches into statistical, classification-based, clustering-based, and information-theoretic methods \cite{ChandolaAnomaly2009}. The isolation forest algorithm, which detects anomalies by measuring random partitions required to isolate data points, has proven effective for identifying outlier behaviour in high-dimensional datasets such as NHI access patterns \cite{LiuIsolationForest2008}.

Research on AI-driven dynamic trust assessment proposes frameworks using both supervised and unsupervised \gls{ml} models to generate adaptive trust scores for identity verification \cite{DynamicTrust2024}. Rather than binary authenticated-or-not decisions, these systems compute continuous trust values reflecting the evolving risk profile of each identity. This approach directly informs risk-scoring engine design for NHI governance. The broader landscape of ML in security operations documents both the promise and challenges of deploying detection systems in production, including scarcity of labelled data and the cost of false positives \cite{MLSecurityOps2023}. The emergence of large language model-based autonomous agents introduces new possibilities for security orchestration with agentic systems capable of reasoning about complex scenarios \cite{AgenticAI2024}.

\section[Policy-as-Code and Governance Automation]{\textbf{Policy-as-Code and Governance Automation}}
The policy-as-code paradigm represents a shift from manual, document-based governance to machine-executable policy definitions that can be version-controlled, tested, and automatically enforced. Systematic reviews identify \gls{opa} with its Rego policy language as the leading framework for expressing governance constraints in declarative, auditable form \cite{PolicyAsCode2023, OPADocs2024}. This enables organisations to codify rules such as maximum credential age and permission scope boundaries as testable code integrating into \gls{cicd} pipelines.

Systematic literature reviews on software secret management document the pervasive challenge of secret sprawl---the uncontrolled proliferation of credentials across repositories and configuration files \cite{SecretSprawl2024}. Studies on automated credential lifecycle management demonstrate that programmatic rotation can reduce exposure windows from months to hours \cite{CredentialRotation2024}. Standards bodies establish rotation schedules and key management practices informing NHI governance policy rules \cite{NIST80057}.

The integration of security into DevOps workflows has been surveyed through practitioner studies. Reports on DevSecOps adoption reveal that automated security testing achieves faster remediation times \cite{GitLabDevSecOps2024}. Performance research demonstrates that elite-performing organisations invest in automated governance controls reducing manual overhead while strengthening security posture \cite{DORAReport2023}. Security challenges specific to microservice architectures, where inter-service credentials grow quadratically with service count, have been systematically analysed \cite{MicroserviceSecurity2023}.

\section[Existing Tools and Limitations]{\textbf{Existing Tools and Their Limitations}}
Several commercial and open-source tools address aspects of the NHI governance problem, though none provides the comprehensive, autonomous governance capability required by modern enterprise environments.

HashiCorp Vault is an industry-standard secrets management platform that provides dynamic secret generation, encryption as a service, and fine-grained access policies based on identity \cite{HashiCorpVault}. It excels as a centralised secrets store and supports programmatic credential rotation through its API. However, Vault does not autonomously discover NHIs across multi-platform environments, does not perform AI-driven risk scoring of credential postures, and does not operate as an agentic governance system capable of independently identifying and remediating policy violations without human intervention.

Microsoft Entra Workload ID provides identity services for workloads running in Azure, supporting managed identities and federated credentials for cloud-native applications \cite{EntraWorkloadID}. It offers basic lifecycle policies such as credential expiry notifications for Azure-native service principals. However, its scope is limited to the Azure ecosystem and it lacks cross-platform NHI correlation with GitHub tokens, \gls{cicd} pipeline secrets, or third-party service credentials. It provides no AI-driven risk classification or anomaly detection capabilities.

Both tools, along with other solutions reviewed, share a common limitation: they address individual components of the NHI governance lifecycle---discovery, storage, rotation, or monitoring---but none integrates all four functions into a continuously operating autonomous agent. The certificate lifecycle management challenges documented in industry research further highlight the operational gaps that emerge when expiry monitoring, policy enforcement, and automated renewal are handled by disconnected tools \cite{Ponemon2024}.

\section[Research Gap and Rationale]{\textbf{Research Gap and Rationale}}
The survey reveals a persistent and well-documented gap in the current landscape of NHI security. Academic research has established the theoretical foundations for zero-trust identity frameworks, AI-driven anomaly detection, and policy-as-code enforcement. Industry reports have quantified the scale and severity of the NHI threat, demonstrating that unmanaged machine credentials represent one of the most significant attack surfaces in modern enterprise environments. Yet no existing solution---whether academic prototype or commercial product---combines autonomous cross-platform NHI discovery, AI-driven continuous risk scoring, policy-as-code enforcement with machine-executable governance rules, and automated credential rotation into a single unified governance agent.

The NHI Governance Agent implemented in this project was designed to close this gap. By integrating LangChain-based agentic AI orchestration for intelligent risk assessment, \gls{opa} Rego policies for codified governance enforcement, and automated rotation via \gls{api}-driven workflows, the system brought together capabilities that had previously existed only in isolation. The agent operates continuously and autonomously, eliminating the manual intervention bottlenecks and inter-tool coordination overhead that characterise current approaches.

\section[Feasibility Study]{\textbf{Feasibility Study}}
The technical feasibility of this project was confirmed through the robust availability and maturity of the required cloud platform APIs, AI orchestration frameworks, and infrastructure automation tools. The Azure \gls{sdk} for Python and the Microsoft Graph \gls{api} provided comprehensive programmatic access to Azure Active Directory service principals, application registrations, and managed identities \cite{AzureSDKPython, MSGraphAPI}. The GitHub \gls{rest} \gls{api} enabled enumeration of deploy keys, Actions secrets, and repository-level NHIs \cite{GitHubRESTAPI}. LangChain provided the agentic orchestration framework used for AI-driven risk scoring and decision-making \cite{LangChainDocs}, while \gls{opa} with Rego supported the declarative policy-as-code enforcement layer \cite{OPADocs2024}. Azure Key Vault offered secure credential storage and rotation capabilities \cite{AzureKeyVault}, and \gls{ga} provided the \gls{cicd} automation platform for executing governance workflows on configurable schedules \cite{GitHubActions2024}. Python~3.11, selected as the primary implementation language, offered mature async capabilities and extensive library support for the required integrations \cite{Python311}. The React~19 framework with TypeScript provided the frontend monitoring dashboard, while scikit-learn supplied the machine learning models for anomaly detection and trend prediction. All technologies were validated as production-ready through their use in the implemented system.

\section*{Chapter Summary}
This chapter surveyed the academic research, industry reports, and commercial tools relevant to NHI governance. The review established that the NHI threat landscape is both severe and rapidly expanding, that zero-trust and AI-driven anomaly detection frameworks provide the theoretical basis for intelligent identity governance, and that policy-as-code approaches enable machine-executable enforcement of governance constraints. Critically, the survey identified a persistent gap: no existing solution integrates cross-platform discovery, AI-driven risk scoring, policy enforcement, and automated rotation into a single autonomous agent. The feasibility study confirmed that the required technologies are mature and available, and their suitability was validated through the implementation described in subsequent chapters. The next chapter presents the system architecture designed to address this gap.
